\documentclass[11pt,a4paper]{amsart}
\usepackage{amsmath,amssymb,amsthm,mathtools}
\usepackage[utf8]{inputenc}
\usepackage{hyperref}
\usepackage{booktabs}

\newcommand{\perm}{\operatorname{perm}}
\newcommand{\rkF}{\operatorname{rk}_{\mathbb{F}_2}}
\newcommand{\Gr}{\operatorname{Gr}}
\newcommand{\CC}{\mathbb{C}}
\newcommand{\FF}{\mathbb{F}}
\newcommand{\ZZ}{\mathbb{Z}}

\newtheorem{theorem}{Theorem}
\newtheorem{lemma}[theorem]{Lemma}
\newtheorem{proposition}[theorem]{Proposition}
\newtheorem{corollary}[theorem]{Corollary}
\theoremstyle{definition}
\newtheorem{definition}[theorem]{Definition}
\newtheorem{remark}[theorem]{Remark}
\newtheorem{conjecture}[theorem]{Conjecture}

\title[Grassmannian Geometry and Modular Separation]{Grassmannian Rank Ceilings and Finite-Characteristic Evaluation Ranks: Permanent versus Determinant}

\author{[Authors]}
\date{\today}

\begin{document}
\begin{abstract}
We study the evaluation matrices of $\det_n$ and $\perm_n$ over finite fields,
focusing on their $\FF_2$-ranks under bipartite row-column partitions.
Via the Cauchy--Binet / Laplace expansion, the two-party evaluation matrix of $\det_n$
explicitly factors through the Plücker coordinates of the row span in the $\{0,1\}$-restricted
Grassmannian $\Gr_{0,1}(k, n, \FF_3)$, yielding an explicit structural rank ceiling.
In contrast, $\perm_n$ admits no such low-dimensional subspace factorization.

We report:
\begin{enumerate}
  \item An explicit bipartite Plücker factorization proving the Grassmannian rank ceiling
        (Lemma~\ref{lem:plucker-factorization}).
  \item Exact projective quotient cardinality: $|P_{0,1}(\FF_3^4)| = 15$ (Theorem~\ref{thm:quotient}).
  \item Exact finite evaluation rank separation at $n = 6$:
        $\rkF([\perm_6 \ne 0 \bmod 3]) = 1456$ versus
        $\rkF([\det_6 \ne 0 \bmod 3]) = 264$ (a $5.52\times$ ratio)
        under the canonical partition (Proposition~\ref{prop:rank-separation}).
  \item A structural Schur $IP_{n-1}$ embedding: on a diagonal slice,
        $\det_n(X) = 1 - IP_{n-1}(u,v) \bmod 3$, identifying PTF lower bound barriers.
  \item Three concrete orbit-LP separating invariants at $n = 4$:
        threshold weight ($182.32$ vs $104.19$), Yao dual support ($127/317$ vs $73/317$),
        and dual-core checkerboard rank ($186$ vs $138$).
\end{enumerate}
\end{abstract}

\maketitle

%% =====================================================================
\section{Bipartite Plücker Factorization and Grassmannian Ceilings}
%% =====================================================================

Let $X \in \{0,1\}^{n \times n}$ be partitioned into two blocks of rows:
$U \in \{0,1\}^{k \times n}$ (the top $k$ rows) and $V \in \{0,1\}^{(n-k) \times n}$ (the bottom $n-k$ rows).
Let $M_{\det_n}(U, V) = [\det(X) \ne 0 \bmod 3]$ denote the bipartite evaluation matrix.

\begin{lemma}[Bipartite Plücker Factorization]\label{lem:plucker-factorization}
Over $\FF_3$, the Laplace expansion along the top $k$ rows yields:
\[
  \det(X) \;=\; \sum_{S \subset [n],\, |S|=k} \operatorname{sgn}(S) \cdot \det(U_{[:, S]}) \cdot \det(V_{[:, S^c]}).
\]
The vector $\mathbf{p}(U) = (\det(U_{[:, S]}))_{|S|=k} \in \FF_3^{\binom{n}{k}}$ represents the Plücker
coordinates of the row space $\operatorname{span}(U) \subset \FF_3^n$.
Because $U$ has $\{0,1\}$-entries, the set of attainable non-zero Plücker vectors modulo scalar
multiplication is bounded by the cardinality of the restricted Grassmannian $|\Gr_{0,1}(k, n, \FF_3)|$.
Consequently, the communication matrix factors through this coordinate map, establishing:
\[
  \rk_{\FF_2} M_{\det_n} \;\le\; |\Gr_{0,1}(k, n, \FF_3)|.
\]
\end{lemma}

\begin{proof}
The evaluation $\det(X) \bmod 3$ is an inner product of the Plücker coordinate vector $\mathbf{p}(U)$
and the complementary minor vector of $V$. If two matrices $U_1, U_2$ have the same row span in $\FF_3^n$,
their Plücker coordinate vectors are proportional, $\mathbf{p}(U_1) = c \cdot \mathbf{p}(U_2)$ for $c \in \FF_3^\times = \{\pm 1\}$.
Hence, the number of distinct row profiles in the evaluation matrix is at most the number of reachable
subspaces in $\Gr_{0,1}(k, n, \FF_3)$, bounding the linear rank over any field.
\end{proof}

\begin{theorem}[Projective Quotient]\label{thm:quotient}
  $|P_{0,1}(\FF_3^4)| = |(\{0,1\}^4 \setminus \{0\}) / \pm 1| = 15$.
\end{theorem}

\begin{proposition}[Evaluation Rank Separation at $n=6$]\label{prop:rank-separation}
Under the canonical bipartite row partition of $\{0,1\}^{6 \times 6}$ with $k = 3$:
\[
  \rkF([\perm_6 \ne 0 \bmod 3]) = 1456, \qquad \rkF([\det_6 \ne 0 \bmod 3]) = 264.
\]
The permanent exhibits a $5.52\times$ rank inflation over the determinant.
\end{proposition}

\begin{remark}[Asymptotic Scope]
We explicitly distinguish this finite-dimensional computational certificate from an asymptotic
lower bound. Extending this separation to super-polynomial complexity lower bounds requires
ruling out arbitrary non-bipartite and adaptive factorizations, which remains an active open challenge.
\end{remark}

%% =====================================================================
\section{The Schur Embedding and PTF Obstructions}
%% =====================================================================

By restricting $\det_n$ to an affine slice with identity blocks and vectors $u, v$,
one obtains the Schur complement identity:
\[
  \det_n \begin{pmatrix} I_{n-1} & u \\ v^T & 1 \end{pmatrix} \;=\; 1 - \langle u, v \rangle \;=\; 1 - IP_{n-1}(u, v).
\]
This embeds the Inner Product modulo~$3$ directly into the determinant.
Consequently, separating the modular determinant from the permanent via polynomial threshold
degree faces standard communication complexity lower bounds of the Inner Product problem.

%% =====================================================================
\section{Orbit-LP Invariants at $n=4$}
%% =====================================================================

Linear programming over the symmetry orbits of the hypercube $\{0,1\}^{4 \times 4}$
yields three explicit separating metrics:
\begin{enumerate}
  \item \textbf{Threshold Weight:} The minimum $\ell_1$-norm of separating hyperplane weights
        is $182.32$ for $\perm_4$ versus $104.19$ for $\det_4$.
  \item \textbf{Yao Dual Support:} The dual witness support cardinality is $127/317$ for $\perm_4$
        versus $73/317$ for $\det_4$.
  \item \textbf{Checkerboard Rank:} The $\FF_2$-rank under the checkerboard partition is
        $186$ for $\perm_4$ versus $138$ for $\det_4$.
\end{enumerate}

\begin{thebibliography}{99}
\bibitem{KN97}
  E.~Kushilevitz and N.~Nisan,
  \emph{Communication Complexity},
  Cambridge University Press, 1997.
\end{thebibliography}

\end{document}
